\section{The delivery evidence layer}

Start with a limit that most engineering organizations already know they have.
The delivery-metrics programs leaders run today, DORA's four keys and DX Core 4,
measure delivery at the team level, but they cannot tell you whether a movement in
those metrics was driven by AI tooling, by process change, or by a quality
trade-off. DX names this the attribution gap directly: standard delivery metrics,
designed before agents wrote code, cannot say what drove the change they measure
\citep{dxcore4}. An organization can watch its change-failure rate move and have
no way to attribute the move to its agents. The 95\% with no measurable P\&L return and the attribution gap are the
same problem seen from two ends: an outcome that cannot be attributed to the work
that produced it.

Closing that gap requires a layer the protocol does not provide, and naming it
correctly matters, because an adjacent discipline already claimed a
similar-sounding term for the opposite job. Context engineering builds a
\emph{context layer} whose purpose is to feed better inputs to agents so they act
well \citep{ibmcontext}. The layer this paper describes runs in the other
direction. Call it the \emph{delivery evidence layer}: its purpose is to feed
outcome evidence to the authorizer (the human or process that makes the trust
decision, later the cadence) so the organization can decide what an agent
is allowed to do next. Context feeds the actor. Evidence feeds the authorizer.
It is one store, and it has four capabilities.

\paragraph{Attribute.} Attribute joins the spec's per-call action records to
version-control and CI/CD events, so a delivered change carries a legible answer
to which autonomous work produced it and through what path. The join is the
implementer's to build. The spec emits attributable, traceable action records
(SEP-414, SEP-2575) and the pipeline emits commits, reviews, and deploys, but
nothing in the protocol stitches the two together. That stitch is the fix for the
attribution gap, and it is the foundational capability the other three depend on.

\paragraph{Evaluate.} Evaluate attaches outcome signals to each agent-authored
change: change-failure rate, revert, rework, override-before-merge, whether the
change held in production, and cost-to-outcome. These are delivery signals, not
activity counts. They answer a different question than how many actions the agent
took. They answer what the actions delivered, and what the delivery cost.
Evaluate is where the conditioning variable of the whole argument, delivered
outcome, gets computed.

\paragraph{Serve.} Serve hands that evidence back as deterministic, typed,
freshness-contracted context, delivered over the spec's own substrate. Here the
release pays a second dividend: full JSON Schema typing (SEP-2106) makes the
served context checkable, required cache scopes and TTLs (SEP-2549) make its
freshness a contract rather than a hope, and the per-call boundary (SEP-2567) is
where it is delivered and re-authorized. The same protocol that emitted the
events carries the answers back, and Serve reaches two consumers over those rails:
the review cadence (defined in the next section) reads the evidence as the
authorizer's input, and the agents
read the distillation as operating context, the current baselines and constraints
they should act within.

\paragraph{Steer.} Steer feeds evaluations into boundary policy: the outcomes a
workflow has delivered set the autonomy it is granted on its next call.
Steer is the capability that turns the join into a loop, and it is where the
protocol earns its place in this argument, because the per-call enforcement point
is the actuator the loop writes to. It is also the capability most easily
overstated, so the next section defines it carefully and states its limits first.

Attribute and Evaluate are the write path; Serve and Steer are the read path.
None of the four is speculative: a join over data an organization already has,
metrics teams already track, a typed read over a governed transport, and policy
at a boundary the spec now provides. What makes them a layer rather than four
disconnected scripts is that they share one store of delivered-outcome evidence,
keyed to the unit an organization wants to govern.
